\section{恢复后缀构造的离线诊断}
\label{sec:construction-diagnostic}

\subsection{固定历史上的构造比较}

为分析依赖映射在训练目标构造中的作用，选取已经保存的全部\LocalConstructionTotal{}个历史训练配对。在相同共同前缀和真实分歧之后，分别执行三种处理：直接沿用参考后缀的原编号；映射依赖后执行并记录结果；在同一映射执行基础上按原规则核验观察与答案。处理方案、调用预算、不可回放记录的登记和案例选取规则均在本次比较前固定。

三种处理均执行既定参考操作。直接拼接在动作被拒绝后继续尝试剩余参考列表，直至提交或达到原24次调用上限，过程中不重新规划。映射执行保留真实工具观察，验收处理再以参考一致性决定样本准入。因此，本节衡量的是后缀构造的可执行性和参考一致性。

本地精确回放得到\LocalConstructionCompared{}个可比较配对，另有\LocalConstructionUnavailable{}个保留为环境观察不一致。后四个均为原历史中被工具拒绝的分歧，动作、错误类别和拒绝状态相同，Python版本产生的\texttt{UnboundLocalError}详情措辞不同。原418分母和全部逐题记录保留，表\ref{tab:construction-diagnostic}以实际可比较的414条执行分支列出结果。

\begin{table}[htbp]
\centering\small
\caption{相同历史下的后缀构造结果（可比较配对414条）}
\label{tab:construction-diagnostic}
\LocalConstructionTable
\end{table}

\subsection{依赖来源与答案一致性的区别}

414个可比较配对的真实分歧均新增了结果句柄。直接拼接存在348条中间依赖错位和391条结束句柄错位，两类可以重叠，合计覆盖全部414条。即使部分操作能够继续执行，引用的结果来源也可能已经偏离参考路径。依赖映射使后缀引用共同成功前缀或新执行的对应结果，414条均完成参考后缀并得到一致答案。

原编号拼接中有9条最终答案文字与参考相同，但其依赖来源仍不一致，执行中均出现工具错误，其中5条还出现内部结果类型或内容差异。这说明构造验收需要同时观察执行路径和答案，单独的答案匹配无法确定是否得到预期参考后缀。对这些情况保留按固定哈希选择的案例，结合句柄来源及完整内部结果解释其差异。

映射执行与映射后验收在本批414条上产生完全相同的候选序列，验收没有新增剔除。重放核验在这里承担质量检查功能；仅关闭该检查未形成不同的训练材料，因此本批诊断支持依赖映射对目标构造的作用，并为后续消融处理的选择提供依据。训练后模型的总体效果由前文2000题评价报告，偏离历史下的续行表现则由下一节同状态评价检验。
